% ---
% filename : "electrotechnique_machines_20260923_217_paires_poles_frequence.tex"
% id : "electrotechnique_machines_20260923_217"
% title : "Paires de pôles et fréquence à 50 Hz d'un moteur asynchrone triphasé"
% domain : "Électrotechnique"
% subdomain : "Machines"
% tags : ["moteur_asynchrone", "triphasé", "fréquence_synchrone", "glissement", "paire_de_pôles"]
% difficulty : 2
% time_solve : 8
% has_octave : true
% octave_files : ["electrotechnique_machines_20260923_217_paires_poles_frequence.m"]
% author : "om"
% ---
\documentclass[a4paper,12pt,answers,addpoints]{exam}

% ---------------------------------------------------------
% LANGUE, ENCODAGE, FONTS
% ---------------------------------------------------------
\usepackage[utf8]{inputenc}

% ---------------------------------------------------------
% GÉOMÉTRIE ET MISE EN PAGE
% ---------------------------------------------------------
\usepackage[left=1.7cm,right=1.5cm,top=2cm,bottom=1cm]{geometry}
\usepackage{microtype}
\usepackage{float}
\usepackage{needspace}

% ---------------------------------------------------------
% MATHÉMATIQUES
% ---------------------------------------------------------
\usepackage{amsmath, amssymb, bm, esvect}

\newcommand{\V}{\overrightarrow}
\def\vect#1{\overrightarrow{\strut#1}}
\providecommand{\Degrees}[0]{\ensuremath{^\circ}}

% ---------------------------------------------------------
% UNITÉS ET GRANDEURS (SIunitx)
% ---------------------------------------------------------
\usepackage{siunitx}
\sisetup{output-decimal-marker={,}}
\DeclareSIUnit \var {var}
\DeclareSIUnit \VA {VA}
\DeclareSIUnit \kVA {kVA}
\DeclareSIUnit[quantity-product={}] \degree{\text{\textdegree}}

% Permet la césure dans les nombres avec unités (évite les Underfull \hbox)
%\AllowBreakAtQQ

% ---------------------------------------------------------
% GRAPHIQUES : TikZ + CircuitikZ (sans tkz-fct qui cause des erreurs spath3)
% ---------------------------------------------------------
\usepackage{tikz}
\usetikzlibrary{
	arrows.meta,
	intersections,
	decorations.pathreplacing,
	%	calligraphy,
	positioning
}

\usepackage[european,straightvoltages,EFvoltages]{circuitikz}
\usepackage{pgfplots}
\pgfplotsset{compat=1.12}

% Symbole étoile-triangle personnalisé
\newcommand{\wye}{%
	\mathbin{\tikz[x=1ex,y=1ex]{%
			\draw[line width=.1ex] (0,0)--(45:1)--++(-45:1) (45:1)--++(0,1);
	}}%
}

% ---------------------------------------------------------
% TABLEAUX
% ---------------------------------------------------------
\usepackage{tabularx}
\usepackage{tabu}
\usepackage{longtable}

% ---------------------------------------------------------
% HYPERLIENS
% ---------------------------------------------------------
\usepackage[colorlinks=true,
menucolor=red,
breaklinks=true,
urlcolor=blue,
linkcolor=blue]{hyperref}

% ---------------------------------------------------------
% COMMANDES PERSONNELLES
% ---------------------------------------------------------
\newcommand\nomauteur{bdminasmorgul@protonmail.com}
\newcommand\dateTe{23.09.2026}
\newcommand\brancheTe{TS}

% ---------------------------------------------------------
% STYLE DE L'EXERCICE
% ---------------------------------------------------------
\pagestyle{headandfoot}
\firstpageheadrule
\runningheadrule

% En-tête avec largeur fixe pour éviter les dépassements
\firstpageheader{\brancheTe\ (\nomauteur)}{Élève : }{\makebox[3.5cm][r]{\dateTe\ \thepage/\numpages}}
\runningheader{\brancheTe\ (\nomauteur)}{Élève : }{\makebox[3.5cm][r]{\dateTe\ \thepage/\numpages}}

% Pied de page
\newcommand{\totalpointtts}{%
	\begin{tikzpicture}[remember picture, overlay]
		\node[anchor=south west, inner sep=0pt, outer sep=0pt]
		at ([xshift=0.5cm,yshift=0.5cm]current page.south west)
		{\rotatebox{90}{Total des points : \numpoints}};
	\end{tikzpicture}%
}

\firstpagefooter{\hspace*{\fill}\totalpointtts}{}{}
\runningfooter{\hspace*{\fill}\totalpointtts}{}{}

% Réglages exam
\printanswers
\addpoints
\pointsinmargin
\bracketedpoints

% ---------------------------------------------------------
% LISTINGS (code)
% ---------------------------------------------------------
\usepackage{xcolor}
\usepackage{listings}

\lstdefinestyle{octavestyle}{
	language=Matlab,
	basicstyle=\ttfamily\small,
	keywordstyle=\color{blue},
	commentstyle=\color{green!50!black},
	stringstyle=\color{red},
	stepnumber=1,
	frame=single,
	breaklines=true,
	breakatwhitespace=true,
	tabsize=4
}

\usepackage{enumitem}
\setlist[enumerate,1]{label=\alph*), ref=\alph*)}
\newenvironment{explication}{%
	\par\noindent\textbf{Explication. }\itshape
}{\par\normalfont}
\renewcommand\partlabel{\arabic{partno}.}
\begin{document}
	\unframedsolutions
	\pointsinmargin
	\bracketedpoints
	
\section*{Préparation TS PQ CFC ELMO,IELE,PELE}
\begin{questions}
\question
D'après sa plaque signalétique, un moteur asynchrone triphasé, fabriqué aux USA, tourne à \qty{850}{[tr/min]} pour une fréquence du réseau de \qty{60}{[Hz]}.
\begin{parts}
\part[3] Combien comporte-t-il de paires de pôles~?
\part[3] Quelle est sa fréquence de rotation s'il est alimenté sous \qty{50}{[Hz]}~?
\end{parts}

\begin{solution}
	
1. Determination du nombre de paires de poles.

La fréquence de rotation d'un moteur asynchrone est toujours legerement inferieure a sa fréquence de synchronisme, donnee par :
\begin{align*}
n_s = \dfrac{60 \cdot f}{p}
\end{align*}
avec $p$ le nombre de paires de poles. Pour $f = \qty{60}{[Hz]}$, les fréquences de synchronisme possibles sont \qty{3600}{[tr/min]} ($p=1$), \qty{1800}{[tr/min]} ($p=2$), \qty{1200}{[tr/min]} ($p=3$), \qty{900}{[tr/min]} ($p=4$), \qty{720}{[tr/min]} ($p=5$), etc. La fréquence de \qty{850}{[tr/min]} n'etant possible que juste en dessous de \qty{900}{[tr/min]} (un moteur asynchrone ne peut tourner qu'en dessous de sa fréquence de synchronisme), on retient :
\begin{align*}
n_s = \qty{900}{[tr/min]} \quad \Rightarrow \quad p = \dfrac{60 \cdot f}{n_s} = \dfrac{60 \cdot 60}{900} = 4
\end{align*}
Le moteur comporte donc 4 paires de poles, soit 8 poles.

2. Calcul du glissement a 60 Hz.

\begin{align*}
s = \dfrac{n_s - n}{n_s} = \dfrac{900 - 850}{900} = 0,0556
\end{align*}
Le glissement vaut environ \qty{5,56}{\%}.

3. fréquence de rotation a 50 Hz.

La nouvelle fréquence de synchronisme vaut :
\begin{align*}
n_{s50} = \dfrac{60 \cdot f}{p} = \dfrac{60 \cdot 50}{4} = \qty{750}{[tr/min]}
\end{align*}
En supposant que le moteur entraine la meme charge, le glissement est identique. Sa fréquence de rotation devient alors :
\begin{align*}
n_{50} = n_{s50} \cdot (1 - s) = 750 \cdot (1 - 0,0556) = \qty{708,3}{[tr/min]}
\end{align*}


\newpage
\begin{verbatim}
% Télécharger OCTAVE depuis https://wiki.octave.org/Using_Octave et l'installer
% Puis copier-coller le code dans OCTAVE
f_usa = 60;          % frequence du reseau americain [Hz]
n     = 850;         % fréquence de rotation relevee [tr/min]

% 1. fréquence de synchronisme la plus proche au-dessus de 850 tr/min a 60 Hz
p = 4;               % nombre de paires de poles recherche
ns_60 = 60 * f_usa / p;
printf("fréquence de synchronisme a 60 Hz : ns = %.0f tr/min\n", ns_60);

% Verification : p recalcule a partir de ns
p_calc = 60 * f_usa / ns_60;
printf("Nombre de paires de poles : p = %.0f (soit %.0f poles)\n", p_calc, 2*p_calc);

% 2. Glissement a 60 Hz
s = (ns_60 - n) / ns_60;
printf("Glissement : s = %.4f soit %.2f %%\n", s, 100*s);

% 3. fréquence a 50 Hz (meme charge => glissement identique)
f_fr = 50;
ns_50 = 60 * f_fr / p_calc;
n_50  = ns_50 * (1 - s);
printf("fréquence de synchronisme a 50 Hz : ns = %.0f tr/min\n", ns_50);
printf("fréquence de rotation a 50 Hz : n = %.1f tr/min\n", n_50);
\end{verbatim}
\end{solution}
\end{questions}
\end{document}